\fontsize{10}{12}\selectfont
\begin{longtable}{L{2.0cm} L{4.0cm} C{4.0cm}}
\caption{Justification for coding instrument categories with examples}\label{tableA1}\\
\toprule
\textbf{Coding categories} &
\textbf{Relationship to visual literature} &
\textbf{Examples} \\
\midrule
\endfirsthead

\toprule
\textbf{Coding categories} &
\textbf{Relationship to visual literature} &
\textbf{Examples (continued)} \\
\midrule
\endhead

\midrule
\multicolumn{3}{r}{\textit{Continued on next page}} \\
\bottomrule
\endfoot

\bottomrule
\endlastfoot

% =========================
% DENOTATIVE
% =========================
\multicolumn{3}{l}{\textbf{Denotative framing elements}}\\
\addlinespace

\textbf{Militants} &
Studies of ISIS and al-Qaeda magazines show the frequent presence of “martyrs” and fighters, with different causes and outcomes for in-group and out-group combatants \parencite{saltman2015a,winkler2018a,winkler2019e,winter2018a}. &
\eximgcap{figures/appendixa/militants_alnaba10.pdf}{\textit{Al-Naba} (10)}
\hfill
\eximgcap{figures/appendixa/militants_almasra_45.pdf}{\textit{Al-Masra} (45)}\\
\addlinespace

\textbf{Death} &
Graphic images of death in wartime highlight the consequences of conflict and who holds the authority to decide who lives or dies \parencite{hoeijer2004a,carlin2012a}. Work on ISIS/AQ magazines shows “about to die” images are more frequent than depictions of dead bodies \parencite{winkler2018a,hanusch2013a}. &
\eximgcap{figures/appendixa/death_dabiq_15.pdf}{\textit{Dabiq} (15)}
\hfill
\eximgcap{figures/appendixa/death_jihadirecollections_1.pdf}{\textit{Jihadi Recollections} (1)}\\
\addlinespace

\textbf{Humans} &
The number of humans in images denotes brotherhood, individual agency, or intended targets \parencite{wilson2017a,winkler2022d}. &
\eximgcap{figures/appendixa/humans_alnaba_23.pdf}{\textit{Al-Naba} (23)}
\hfill
\eximgcap{figures/appendixa/humans_inspire_3.pdf}{\textit{Inspire} (3)}\\
\addlinespace

\textbf{Destruction} &
Explosions and destruction of infrastructure or religious iconography are used to instill fear and acquiescence \parencite{winkler2018a,winkler2019e}. &
\eximgcap{figures/appendixa/destruction_dabiq_11.pdf}{\textit{Dabiq} (11)}
\hfill
\eximgcap{figures/appendixa/destruction_almasra_4.pdf}{\textit{Al-Masra} (4)}\\
\addlinespace

\textbf{Leaders} &
Images of leaders denote enemies, potential allies, and intra-jihadist conflict following ISIS’s break from al-Qaeda \parencite{winkler2022d}. &
\eximgcap{figures/appendixa/leaders_rumiyah_1.pdf}{\textit{Rumiyah} (1)}
\hfill
\eximgcap{figures/appendixa/leaders_jihadirecollection_4.pdf}{\textit{Jihadi Recollections} (4)}\\
\addlinespace

\textbf{Flags} &
Flags are used to recruit supporters and anchor enmity toward rival states and groups \parencite{el2021a,warrick2016a}. &
\eximgcap{figures/appendixa/flags_dabiq_13.pdf}{\textit{Dabiq} (13)}
\hfill
\eximgcap{figures/appendixa/flags_almasra_12.pdf}{\textit{Al-Masra} (12)}\\
\addlinespace

% =========================
% SEMIOTIC
% =========================
\multicolumn{3}{l}{\textbf{Semiotic framing elements}}\\
\addlinespace

\textbf{Viewer position} &
Camera angle shapes perceived strength and credibility \parencite{kraft1986a,mccain1977a}. &
\eximgcap{figures/appendixa/viewerposition_alnaba_123.pdf}{\textit{Al-Naba} (123)}
\hfill
\eximgcap{figures/appendixa/viewerposition_inspire_2.pdf}{\textit{Inspire} (2)}\\
\addlinespace

\textbf{Image position} &
Foregrounding increases salience while backgrounding provides contextual grounding \parencite{stone2003b}. &
\eximgcap{figures/appendixa/imageposition_dabiq_10.pdf}{\textit{Dabiq} (10)}
\hfill
\eximgcap{figures/appendixa/imageposition_almasra_14.pdf}{\textit{Al-Masra} (14)}\\
\addlinespace

\textbf{Viewer distance} &
Intimate versus public distance implies relational closeness or collective anonymity \parencite{jewitt2008a}. &
\eximgcap{figures/appendixa/viewerdistance_rumiyah_12.pdf}{\textit{Rumiyah} (12)}
\hfill
\eximgcap{figures/appendixa/viewerdistance_inspire_11.pdf}{\textit{Inspire} (11)}\\
\addlinespace

\textbf{Eye contact} &
Direct gaze communicates dominance while averted gaze suggests reduced confrontation \parencite{tang2015a}. &
\eximgcap{figures/appendixa/eyecontact_dabiq_6.pdf}{\textit{Dabiq} (6)}
\hfill
\eximgcap{figures/appendixa/eyecontact_jihadirecollections_1.pdf}{\textit{Jihadi Recollections} (1)}\\
\addlinespace

\textbf{Facial expressions} &
Positive expressions promote trust whereas negative expressions prompt skepticism \parencite{forgas2008a}. &
\eximgcap{figures/appendixa/facialexpression_alnaba_16.pdf}{\textit{Al-Naba} (16)}
\hfill
\eximgcap{figures/appendixa/facialexpression_inspire_4.pdf}{\textit{Inspire} (4)}\\
\addlinespace

\textbf{Stance} &
Body posture signals submissiveness or strength and control \parencite{green2015a,jewitt2008a}. &
\eximgcap{figures/appendixa/stance_rumiyah_10.pdf}{\textit{Rumiyah} (10)}
\hfill
\eximgcap{figures/appendixa/stance_inspire_15.pdf}{\textit{Inspire} (15)}\\
\addlinespace

% =========================
% CONNOTATIVE
% =========================
\multicolumn{3}{l}{\textbf{Connotative framing elements}}\\
\addlinespace

\textbf{State-building} &
Images of governance, services, markets, and territory signal institutional capacity and caliphate aspirations \parencite{damanhoury2022a,el2021a,lokmanoglu2020a,winkler2019a,zelin2015a}. &
\eximgcap{figures/appendixa/statebuilding_alnaba_2.pdf}{\textit{Al-Naba} (2)}
\hfill
\eximgcap{figures/appendixa/statebuilding_almasra_7.pdf}{\textit{Al-Masra} (7)}\\
\addlinespace

\textbf{Law enforcement} &
Images of capturing and punishing lawbreakers reinforce deterrence and community policing narratives \parencite{barr2017a,chouliaraki2018a,damanhoury2018a}. &
\eximgcap{figures/appendixa/lawenforcement_dabiq_4.pdf}{\textit{Dabiq} (4)}
\hfill
\eximgcap{figures/appendixa/lawenforcement_almasra_17.pdf}{\textit{Al-Masra} (17)}\\
\addlinespace

\textbf{Allegiance pledges} &
Allegiance imagery visualizes coalition-building and collective solidarity \parencite{gregg2023a}. &
\eximgcap{figures/appendixa/allegiencepledges_alnaba_12.pdf}{\textit{Al-Naba} (12)}
\hfill
\eximgcap{figures/appendixa/allegiencepledges_almasra_36.pdf}{\textit{Al-Masra} (36)}\\
\addlinespace

\textbf{Media propaganda} &
Infographics and posters function as authority strategies and publicity for media products \parencite{glausch2020a,abdelrahim2019a}. &
\eximgcap{figures/appendixa/mediapropaganda_dabiq_13.pdf}{\textit{Dabiq} (13)}
\hfill
\eximgcap{figures/appendixa/mediapropaganda_inspire_11.pdf}{\textit{Inspire} (11)}\\
\addlinespace

\textbf{About to die} &
“About to die” images are especially circulation-prone and vary by certainty and probability \parencite{zelizer2010a,winkler2016a}. &
\eximgcap{figures/appendixa/abouttodie_rumiyah_3.pdf}{\textit{Rumiyah} (3)}
\hfill
\eximgcap{figures/appendixa/abouttodie_inspire_16.pdf}{\textit{Inspire} (16)}\\
\addlinespace

\textbf{Religion} &
Religious gestures and iconography signal piety, legitimacy, and ideological commitment \parencite{winkler2022a,heck2017a}. &
\eximgcap{figures/appendixa/religion_dabiq_7.pdf}{\textit{Dabiq} (7)}
\hfill
\eximgcap{figures/appendixa/religion_jihadirecollections_2.pdf}{\textit{Jihadi Recollections} (2)}\\
\end{longtable}
